% TODO make hw stylesheet
\documentclass[11pt]{scrartcl}
\usepackage[colorlinks]{hyperref}
\usepackage[a4paper, total={6in, 8in}]{geometry}
\usepackage[usenames,dvipsnames,svgnames]{xcolor}
\usepackage[shortlabels]{enumitem}
\usepackage[framemethod=TikZ]{mdframed}
\usepackage{amsmath,amssymb,amsthm}
\usepackage{multicol}
\usepackage{tabularx}
\usepackage{stmaryrd}

\hypersetup{
	colorlinks=true,
	linkcolor=blue,
	filecolor=magenta,
	urlcolor=magenta,
	pdftitle={MAT 324 Notes},
	hypertexnames=false
}

\usepackage{microtype}
\usepackage{mathtools}
\usepackage[headsepline]{scrlayer-scrpage}
\usepackage{thmtools}
\usepackage[textsize=footnotesize]{todonotes}

\mdfdefinestyle{mdbluebox}{roundcorner=10pt,innerbottommargin=9pt,
	linecolor=blue,backgroundcolor=TealBlue!5,}
\declaretheoremstyle[headfont=\sffamily\bfseries\color{MidnightBlue},
mdframed={style=mdbluebox},]{thmbluebox}

\mdfdefinestyle{mdbloobox}{frametitlefont=\bfseries,innerbottommargin=8pt,
	nobreak=true,backgroundcolor=TealBlue!5,linecolor=blue,}
\declaretheoremstyle[headfont=\bfseries\color{MidnightBlue},
mdframed={style=mdbloobox},headpunct={\\[3pt]},postheadspace=0pt,]{thmbloobox}

\mdfdefinestyle{mdredbox}{frametitlefont=\bfseries,innerbottommargin=8pt,
	nobreak=true,backgroundcolor=Salmon!5,linecolor=RawSienna,}
\declaretheoremstyle[headfont=\bfseries\color{RawSienna},
mdframed={style=mdredbox},headpunct={\\[3pt]},postheadspace=0pt,]{thmredbox}

\mdfdefinestyle{mdgreenbox}{linecolor=ForestGreen,backgroundcolor=ForestGreen!5,
	linewidth=2pt,rightline=false,leftline=true,topline=false,bottomline=false,}
\declaretheoremstyle[headfont=\bfseries\sffamily\color{ForestGreen!70!black},
mdframed={style=mdgreenbox},headpunct={ --- },]{thmgreenbox}

\mdfdefinestyle{mdorangebox}{linecolor=RedOrange,backgroundcolor=RedOrange!5,
	linewidth=2pt,rightline=false,leftline=true,topline=false,bottomline=false,}
\declaretheoremstyle[headfont=\bfseries\sffamily\color{RedOrange!70!black},
mdframed={style=mdorangebox},headpunct={ --- },]{thmorangebox}

\mdfdefinestyle{mdblackbox}{linecolor=black,backgroundcolor=RedViolet!5!gray!5,
	linewidth=3pt,rightline=false,leftline=true,topline=false,bottomline=false,}
\declaretheoremstyle[mdframed={style=mdblackbox}]{thmblackbox}

\declaretheoremstyle[spaceabove=6pt,spacebelow=6pt]{basehead}

\declaretheorem[style=thmbluebox,name=Theorem,numberwithin=section]{theorem}
\declaretheorem[style=thmredbox,name=Example,sibling=theorem]{example}
\declaretheorem[style=thmbluebox,name=Corollary,sibling=theorem]{corollary}
\declaretheorem[style=thmbluebox,name=Proposition,sibling=theorem]{prop}
\declaretheorem[style=thmbluebox,name=Lemma,sibling=theorem]{lemma}
\declaretheorem[style=thmbluebox,name=Lemma,numberwithin=theorem]{thmlemma}
\declaretheorem[style=thmbluebox,name=Theorem,numbered=no]{theorem*}
\declaretheorem[style=thmbluebox,name=Lemma,numbered=no]{lemma*}
\declaretheorem[style=thmgreenbox,name=Claim,numberwithin=theorem]{claim}
\declaretheorem[style=thmgreenbox,name=Claim,numbered=no]{claim*}
\declaretheorem[style=thmblackbox,name=Remark,numberwithin=theorem]{remark}
\declaretheorem[style=thmblackbox,name=Remark,numbered=no]{remark*}
\declaretheorem[style=thmgreenbox,name=Definition,sibling=theorem]{definition}
\declaretheorem[style=thmgreenbox,name=Definition,numbered=no]{definition*}
\declaretheorem[style=thmorangebox,name=Warning,sibling=theorem]{warning}
\declaretheorem[style=thmorangebox,name=Warning,numbered=no]{warning*}
\declaretheorem[style=thmblackbox,name=Exercise,sibling=theorem]{exercise}
\declaretheorem[style=thmblackbox,name=Exercise,numbered=no]{exercise*}

\def\equationautorefname~#1\null{(#1)\null}

\addtolength{\textheight}{3.14cm}
\providecommand{\ol}{\overline}
\providecommand{\ceil}[1]{\left\lceil #1 \right\rceil}
\providecommand{\floor}[1]{\left\lfloor #1 \right\rfloor}
\newcommand{\dd}{\mathrm{d}}
\providecommand{\ul}{\underline}
\providecommand{\eps}{\varepsilon}
\providecommand{\half}{\frac{1}{2}}
\providecommand{\dang}{\measuredangle} %% Directed angle
\providecommand{\CC}{\mathbb C}
\providecommand{\FF}{\mathbb F}
\providecommand{\NN}{\mathbb N}
\providecommand{\QQ}{\mathbb Q}
\providecommand{\RR}{\mathbb R}
\providecommand{\ZZ}{\mathbb Z}
\providecommand{\dg}{^\circ}
\providecommand{\ii}{\item}
\providecommand{\setdiff}{\smallsetminus}
\providecommand{\alert}[1]{{\sffamily\textbf{\textcolor{blue}{#1}}}}
\providecommand{\inv}{^{-1}}
\providecommand{\mb}[1]{\mathbf{#1}}

\reversemarginpar

\newenvironment{soln}{\begin{proof}[Solution]}{\end{proof}}
\newenvironment{walkthrough}{\noindent \textbf{Walkthrough.}}{\medskip}
\newlist{walk}{enumerate}{3}
\setlist[walk]{label=\bfseries (\alph*)}

\providecommand{\pow}{\operatorname{Pow}}
\providecommand{\vocab}[1]{{\textbf{\textcolor{ForestGreen}{#1}}}}
\providecommand{\lcm}{\operatorname{lcm}}
\providecommand{\abs}[1]{\left\lvert #1\right\rvert}

\usepackage{asymptote}
\begin{asydef}
	defaultpen(fontsize(10pt));
	unitsize(1cm);
	size(8cm); // set a reasonable default
	usepackage("amsmath");
	usepackage("amssymb");
	settings.tex="pdflatex";
	settings.outformat="pdf";
	// Replacement for olympiad+cse5 which is not standard
	import geometry;
	// recalibrate fill and filldraw for conics
	void filldraw(picture pic = currentpicture, conic g, pen fillpen=defaultpen, pen drawpen=defaultpen)
	{ filldraw(pic, (path) g, fillpen, drawpen); }
	void fill(picture pic = currentpicture, conic g, pen p=defaultpen)
	{ filldraw(pic, (path) g, p); }
	// some geometry
	pair foot(pair P, pair A, pair B) { return foot(triangle(A,B,P).VC); }
	pair orthocenter(pair A, pair B, pair C) { return orthocentercenter(A,B,C); }
	pair centroid(pair A, pair B, pair C) { return (A+B+C)/3; }
	// cse5 abbreviations
	path CP(pair P, pair A) { return circle(P, abs(A-P)); }
	path CR(pair P, real r) { return circle(P, r); }
	pair IP(path p, path q) { return intersectionpoints(p,q)[0]; }
	pair OP(path p, path q) { return intersectionpoints(p,q)[1]; }
	path Line(pair A, pair B, real a=0.6, real b=a) { return (a*(A-B)+A)--(b*(B-A)+B); }
	// cse5 more useful functions
	picture CC() {
		picture p=rotate(0)*currentpicture;
		currentpicture.erase();
		return p;
	}
	pair MP(Label s, pair A, pair B = plain.S, pen p = defaultpen) {
		Label L = s;
		L.s = "$"+s.s+"$";
		label(L, A, B, p);
		return A;
	}
	pair Drawing(Label s = "", pair A, pair B = plain.S, pen p = defaultpen) {
		dot(MP(s, A, B, p), p);
		return A;
	}
	path Drawing(path g, pen p = defaultpen, arrowbar ar = None) {
		draw(g, p, ar);
		return g;
	}
\end{asydef}

\usepackage{mathrsfs}

\ihead{\footnotesize MAT 324 Notes}
\ohead{\footnotesize \today}

\title{\vspace{-2em}MAT 324: Real Analysis}
\author{\large Aditya Pahuja}
\date{}
\begin{document}
\maketitle
\tableofcontents

\pagebreak

\section{August 26}\label{sec:aug26}

The idea of measure is to measure the ``size'' of subsets of $\RR$
(or of other spaces).

\begin{exercise}
    [Think about briefly before reading on]
    What are some reasonable properties that a size function
    $\abs{\bullet}\colon 2^\RR\to\RR_{\ge0}\cup\{\infty\}$ should satisfy?
    (Henceforth in these notes, we will denote
    $\RR_{\ge0}\cup\{\infty\}$ by $[0,\infty]$.)
\end{exercise}

The four properties that we will start by asking for
in a size function on $\RR$ are
\begin{enumerate}[(1)]
    \setcounter{enumi}{-1}
    \ii Intervals behave like they're supposed to:
    if $I\subseteq\RR$ is an interval with endpoints $a\le b$,
    then $|I| = b-a$.
    \ii The size function preserves order:
    if $A\subseteq B\subseteq\RR$, then $|A|\le |B|$.
    \ii The size function satisfies countable subadditivity:
    if $A_1,A_2,\ldots\subseteq\RR$ is a sequence of subsets of $\RR$, then
    \begin{equation*}
	\abs{\bigcup_{n=1}^\infty A_n} \le \sum_{n=1}^\infty \abs{A_n}.
    \end{equation*}
    This is the least natural one, in my opinion;
    in particular, it's more natural to ask instead for
    some form of additivity, like
    \begin{equation*}
	|A\cup B| = |A| + |B|
    \end{equation*}
    when $A$ and $B$ are disjoint.
    However, it turns out this is too much to ask for
    in conjunction with the properties in this list;
    see \autoref{prop:additivity-fails}.
    We ask about countably many sets instead of finitely many
    because analysis is all about limiting behavior,
    which is of course intimately related to infinite sums.
    \ii Translation preserves size:
    if $A\subseteq\RR$ and $t\in\RR$, then, defining
    $t + A = \{t + a : a\in A\}$, we want $\abs{t+A} = \abs{A}$.
\end{enumerate}

A natural notion of such a size function arises
from the idea of trying to cover $A\subseteq\RR$ with
a countable collection of intervals and measuring their length:
\begin{definition}[Outer measure]
    \label{def:outermeasure}
    Given interval $I$ with endpoints $a\le b$,
    define its length $\ell(I)$ to be $b-a$.
    Then, the \vocab{outer measure} of a subset $A\subseteq\RR$ is
    \begin{equation*}
        |A| \coloneq
	\inf\left\{
	    \sum_{n=1}^\infty \ell(I_n) :
	    I_n\text{ are open intervals s.t. }
	    A\subseteq \bigcup_{n=1}^\infty I_n
	\right\}.
    \end{equation*}
    Let $\mathcal L_A$ be the set whose infimum is being taken.
\end{definition}

\begin{lemma}[HW01 Problem 4]
    \label{lem:om-endpoints-irrelevant}
    In the definition of $\mathcal L_A$,
    one can freely remove the word ``open''
    or replace it with ``closed'' or ``half-open''
    without altering the resulting function $\abs{\bullet}\colon 2^\RR\to[0,\infty]$.
\end{lemma}

\begin{corollary}[From the 08/28 lecture]
    \label{coro:om-countable-zero}
    The outer measure of a countable set is zero.
\end{corollary}

\begin{proof}
    The stupidest way to do this:
    let $A = \{a_1, a_2, \dots\}$ be countable; then,
    the sequence of intervals $I_n = [a_n, a_n] = \{a_n\}$ covers $A$, so
    \begin{equation*}
        0\le\abs{A}\le \sum_{n=1}^\infty \ell(I_n) = a_n-a_n = 0.\qedhere
    \end{equation*}
\end{proof}

Now we show that our outer measure is a satisfactory size function
in the sense of the desiderata detailed above.

\begin{prop}
    \label{prop:om-properties}
    The outer measure satisfies properties (0)--(3).
\end{prop}

\begin{proof}
    (0) Let $I$ be an interval. Using \autoref{lem:om-endpoints-irrelevant},
    we immediately deduce that $|I|\le\ell(I)$,
    since $I$ can be covered by the sequence $I$, $\varnothing$, $\varnothing$, \dots.

    For the opposite inequality, we consider several cases.
    If $I$ is a closed, bounded interval then for a sequence of
    bounded open intervals $I_1$, $I_2$, \dots covering $I$,
    we can use compactness to restrict to finitely many of these intervals
    and ultimately show that
    \begin{equation*}
	\sum_{n=1}^\infty \ell(I_n) \ge \ell(I)
    \end{equation*}
    (see the solution to Problem 8a on the first day quiz).

    If $I$ is a generic bounded interval with endpoints $a$ and $b$ then
    \begin{equation*}
	b-a+2\eps = \abs{[a+\eps,b-\eps]} \le \abs I
    \end{equation*}
    for all $\eps > 0$, using monotonicity (property (1)),
    so $\abs I = b-a = \ell(I)$.

    Finally, if $I$ is unbounded, then any sequence of intervals covering $I$
    must cover an interval of length $n$ for all positive integers $n$;
    for example, if $I = (0,\infty)$, then a cover of $I$ by intervals
    must also cover $(0,n)$ for all $n$. Thus, by the logic for the previous cases,
    the total length of all the intervals in the cover must be at least $n$
    for all $n$, i.e. the total length is $\infty$, so $\abs I = \ell(I) = \infty$.

    (1) If $A\subseteq B$ then $\mathcal L_B\subseteq\mathcal L_A$
    because any sequence of intervals that covers $B$ also covers $A$, so
    \begin{equation*}
        |A| = \inf\mathcal L_A \le \inf\mathcal L_B = |B|.
    \end{equation*}

    (2) Let $A_1$, $A_2$, \dots be a sequence of subsets of $\RR$.
    By definition of infimum, for each $\eps > 0$
    and for each $n$ there is a sequence of intervals
    $I^{(n)}_1$, $I^{(n)}_2$, \dots such that
    \begin{equation*}
	A_n\subseteq \bigcup_{m=1}^\infty I^{(n)}_m,\quad
	\abs{A_n} + \frac{\eps}{2^n} \ge \sum_{m=1}^\infty \ell(I^{(n)}_m).
    \end{equation*}
    The union condition also means that $\bigcup_{n=1}^\infty A_n$
    is contained in the union of the $I_m^{(n)}$, of which there are countably many
    because $\NN\times\NN$ is countable, so
    using the above inequality and the definition of outer measure again,
    \begin{equation*}
	\sum_{n=1}^\infty \left(\abs{A_n} + \frac\eps{2^n}\right)
	= \left(\sum_{n=1}^\infty \abs{A_n}\right) + \eps
	\ge \sum_{n=1}^\infty \sum_{m=1}^\infty \ell(I_m^{(n)})
	\ge \abs{\bigcup_{n=1}^\infty A_n}.
    \end{equation*}
    This holds for all $\eps>0$, so
    $\sum_{n=1}^\infty\abs{A_n}\ge\abs{\bigcup_{n=1}^\infty A_n}$.

    (3) If $I_1$, $I_2$, \dots is a sequence of intervals covering $A$,
    then $t+I_2$, $t+I_2$, \dots covers $t+A$, so
    \begin{equation*}
	\mathcal L_{t+A}\subseteq\mathcal L_A\implies |A|\le |t+A|.
    \end{equation*}
    The converse follows from $|t+A|\le \lvert-t + (t+A)\rvert = |A|$
    because $-t + (t+A) = A$.
\end{proof}

\section{August 28}\label{sec:aug28}

It would be nice if we had
\begin{equation}\tag{\textasteriskcentered}
    \label{eq:om-finite-additivity}
    |A\cup B| = |A| + |B|
\end{equation}
whenever $A\cap B=\varnothing$; we shall explore the extent to which this holds,
in terms of only the properties (0)--(3) (as opposed to appealing to
the interval-covering definition).

First, let's justify that additivity is a very strong and useful property.

\begin{exercise}
    Show that \autoref{eq:om-finite-additivity} implies (1)
    and the finite version of (2), i.e. $\abs{A\cup B}\le \abs{A}+\abs{B}$.
    Show also that it implies that $\abs{\varnothing}$ is $0$
    unless the size of every subset of $\RR$ is $\infty$.
\end{exercise}

In fact, \autoref{eq:om-finite-additivity} lets us recover countable additivity.
\begin{prop}
    \label{prop:finadd-and-ctablesubadd-implies-ctableadd}
    If $\abs\bullet\colon 2^\RR\to[0,\infty]$ is
    such that \autoref{eq:om-finite-additivity} and property (2) hold, then
    \begin{equation*}
	\abs{\bigcup_{n=1}^\infty A_n} = \sum_{n=1}^\infty \abs{A_n}
    \end{equation*}
    for any sequence $A_1$, $A_2$, \dots of subsets of $\RR$
    such that $A_i\cap A_j=\varnothing$ whenever $i\ne j$.
\end{prop}

\begin{proof}
    Since the left-hand side is less than the right-hand side by (2),
    we only need to prove the reverse inequality.
    This is fairly quick: we have
    \begin{equation*}
	\abs{\bigcup_{n=1}^\infty A_n} \overset{(1)}\ge
	\abs{\bigcup_{n=1}^N A_n} \overset{(*)}= \sum_{n=1}^N \abs{A_n}
    \end{equation*}
    (noting that \autoref{eq:om-finite-additivity} implies monotonicity),
    and taking the limit as $N$ goes to infinity gives
    \begin{equation*}
	\abs{\bigcup_{n=1}^\infty A_n} \ge \sum_{n=1}^N \abs{A_n}.\qedhere
    \end{equation*}
\end{proof}

This would be an incredible property to have,
as we have upgraded countable subadditivity drastically.

Let us now prove results in the other direction
that clue us to the extent to which additivity holds in reality.
\begin{prop}
    \label{prop:om-zero-meas-addend}
    Suppose $\abs\bullet\colon 2^\RR\to[0,\infty]$ is such that
    (1) and (2) hold.
    Then, if $A$ and $B$ are disjoint subsets of $\RR$
    such that $\abs A = 0$ or $\abs B = 0$, we have
    \begin{equation*}
	\abs{A\cup B} = \abs A+\abs B.
    \end{equation*}
\end{prop}

\begin{proof}
    Suppose without loss of generality that $\abs B = 0$.
    Then, the desired equation follows from the inequalities
    \begin{equation*}
	\abs A+\abs B \overset{(2)}\ge \abs{A\cup B}
	\overset{(1)}\ge \abs A
	= \abs A + \abs B.\qedhere
    \end{equation*}
\end{proof}

This is a very basic case; in fact, the $A\cap B = \varnothing$
condition is redundant.
The next one is much more substantive:
\begin{prop}
    \label{prop:om-open-closed-addend}
    Suppose $\abs\bullet\colon2^\RR\to[0,\infty]$ is the outer measure on $\RR$.
    If $A$ and $B$ are disjoint subsets of $\RR$
    such that either of $A$ and $B$ are open or closed, then
    \begin{equation*}
	\abs{A\cup B} = \abs A + \abs B.
    \end{equation*}
\end{prop}

\begin{proof}
    As has been done earlier, we only need to prove that
    the left-hand side is at least as large as the right-hand side,
    since the other direction is immediate from (2).

    We will prove the case where (WLOG) $B$ is open first.
    To start with, suppose that $B$ is an open interval $(a,b)$.
    If $I_1$, $I_2$, \dots is a sequence of intervals
    that covers $A\cup B$, define
    \begin{align*}
	I_k^- &\coloneq (-\infty, a]\cap I_k\\
	I_k^0 &\coloneq (a,b)\cap I_k = B\cap I_k\\
	I_k^+ &\coloneq [b, \infty)\cap I_k.
    \end{align*}
    Since $A$ lives in $\RR\setminus B$, we see that
    $A$ is contained in $\bigcup_{k=1}^\infty I_k^-\cup\bigcup_{k=1}^\infty I_k^+$,
    while $B$ is contained in (in fact, equal to) $\bigcup_{k=1}^\infty I_k^0$.
    Thus,
    \begin{equation*}
	\abs A + \abs B \le \sum_{k=1}^\infty \ell(I_k^-) +
	\sum_{k=1}^\infty \ell(I_k^+) + \sum_{k=1}^\infty\ell (I_k^0)
	= \sum_{k=1}^\infty (\ell(I_k^-) + \ell(I_k^0) + \ell(I_k^+))
	= \sum_{k=1}^\infty \ell(I_k).
    \end{equation*}
    Taking the infimum over all sequences of intervals that cover $A\cup B$,
    we deduce that $\abs A + \abs B \le \abs{A\cup B}$ as desired.
    This can be extended to the case where $B$ is a union of finitely many intervals
    by a simple induction argument on the number of intervals in $B$.

    To extend this to arbitrary open sets, we first prove the following:
    \begin{claim}
        \label{claim:open-set-into-intervals}
        An open set $U\subseteq\RR$ can be written as
	the union of countably many pairwise disjoint open intervals.
    \end{claim}
    \begin{proof}
        The connected components of $U$ are open in $U$,
	hence open in $\RR$, so they must be open intervals,
	and of course they are pairwise disjoint
	because they are the connected components of $U$.
	To show that $U$ has countably many connected components,
	we note that each connected component of $U$,
	being an open interval, contains a rational number.
	Consider a function sending each connected component
	to a rational number within the component;
	this is an injective function that maps the components into $\QQ$,
	so the set of components must have cardinality
	at most that of $\QQ$, hence is countable.
    \end{proof}

    Write $B$ as the union of countably many open intervals $I_1$, $I_2$, \dots.
    For a positive integer $N$, we can use monotonicity
    and the case where $B$ is a union of finitely many intervals to get
    \begin{equation*}
	\abs{A\cup B} = \abs{A\cup \bigcup_{n=1}^\infty I_n} \overset{(1)}\ge
	\abs{A\cup \bigcup_{n=1}^N I_n} = \abs A + \sum_{n=1}^N \abs{I_n}.
    \end{equation*}
    Then, taking the limit as $N$ grows arbitrarily large gives
    \begin{equation*}
	\abs{A\cup B} \ge \abs A + \sum_{n=1}^\infty \abs{I_n}
	\overset{(2)}\ge \abs A + \abs B.
    \end{equation*}

    Lastly, suppose $B$ is closed. Then, if $A\cup B$ is covered by
    open intervals $I_1$, $I_2$, \dots, let $U = \bigcup_{n=1}^\infty I_n$.
    Clearly, $U\setminus B$ is open, so $\abs U = \abs B + \abs{U\setminus B}$
    since we have proven additivity when one of the sets is open. Then
    \begin{equation*}
	\sum_{n=1}^\infty \ell(I_n) \overset{(2)}\ge
	\abs U = \abs B + \abs{U\setminus B} \overset{(1)}\ge \abs B + \abs A.
    \end{equation*}
    Taking the infimum over all coverings of $A\cup B$ by intervals
    gives $\abs{A\cup B}\ge \abs B + \abs A$.
\end{proof}

This is pretty great; our size function is additive
on a large collection of pretty reasonable sets.
As alluded to, though, additivity does not hold in general:
some sets are so incredibly messy that they do not allow for additivity.
\begin{prop}
    \label{prop:additivity-fails}
    Let $\abs\bullet\colon2^\RR\to[0,\infty]$
    satisfy properties (0)--(3).
    There exist disjoint subsets $A$, $B$ of $\RR$
    such that $\abs{A\cup B} < \abs A + \abs B$.
\end{prop}

\begin{proof}
    Define the equivalence relation $\sim$ on $[-1,1]$ by
    $x\sim y$ if and only if $x-y\in\QQ$ (check that it really is one!),
    and let $[x]$ (which is a subset of $[-1,1]$)
    be the equivalence class of $x$ with respect to $\sim$.
    The centerpiece of our construction will be a set $S$
    defined by adding to $S$ exactly one element of each equivalence class;
    that is, $\#(S\cap [x]) = 1$ for all $x\in[-1,1]$.
    Since $S$ has exactly one representative of each equivalence class,
    if $q$ and $r$ are two distinct rational numbers
    then $q+S$ and $r+S$ must be disjoint, since otherwise
    $q+x = r+y \iff x-y = q-r \iff x\sim y$ for some distinct $x,y\in S$.

    Note now that if $x\sim y$ for $x,y\in[-1,1]$, then $x-y\in[-2,2]$.
    Letting $q_1$, $q_2$, \dots be an enumeration of
    the countable set $[-2,2]\cap\QQ$, we thus have
    \begin{equation*}
	[-1,1]\subseteq\bigcup_{n=1}^\infty (q_n + S)
	\implies \abs{[-1,1]} \overset{(0)} = 2
	\overset{(1)}\le \abs{\bigcup_{n=1}^\infty (q_n + S)}
	\overset{(2)}\le \sum_{n=1}^\infty\abs{q_n+S}
	\overset{(3)}= \sum_{n=1}^\infty\abs{S}.
    \end{equation*}
    This implies that $\abs S\ne 0$, as otherwise the right-hand side
    of the inequality would be zero.

    On the other hand, the union $\bigcup_{n=1}^\infty (q_n + S)$
    is contained in $[-3,3]$ because $q_n\in[-2,2]$
    and each element of $S$ is in $[-1,1]$, so
    \begin{equation*}
	\abs{\bigcup_{n=1}^\infty (q_n + S)} \overset{(1)}\le
	\abs{[-3,3]} \overset{(0)}= 6.
    \end{equation*}
    Since $\abs{\bigcup_{n=1}^N (q_n + S)}$ is bounded above by $6$ for all $N$
    while $\sum_{n=1}^N \abs{q_n+S} = N\cdot\abs S$
    grows arbitrarily large with $N$, there is some minimal $N_0$ for which
    \begin{equation*}
        \abs{\bigcup_{n=1}^{N_0} (q_n + S)} < N_0\cdot \abs S.
    \end{equation*}
    Of course $N_0 > 1$ (since equality holds trivially at $N=1$),
    so by minimality of $N_0$, we have
    \begin{equation*}
	\abs{\bigcup_{n=1}^{N_0-1}(q_n + S)} = (N_0-1)\cdot\abs S
    \end{equation*}
    which implies that
    \begin{equation*}
	\abs{(q_{N_0} + S)\cup \bigcup_{n=1}^{N_0-1}(q_n + S)}
	< N_0\cdot\abs S =
	\abs{q_{N_0} + S} + \abs{\bigcup_{n=1}^{N_0-1} (q_N + S)},
    \end{equation*}
    so $q_{N_0} + S$ and $\bigcup_{n=1}^{N_0-1}(q_N+S)$ are
    the desired $A$ and $B$.
\end{proof}

In English, what happened was that if additivity were to hold, then
$\abs{\bigcup_{n=1}^\infty (q_n+S)}$ should be either zero or $\infty$;
but our construction somehow forced a middle ground because
the translates of $S$ are somehow ``too interlocked with each other'';
$S$ is simultaneously too big to have a size of zero
while too small for its infinitely many disjoint translates
to have infinite size together.

\section{September 2}\label{sec:sep2}
As we saw last time, some sets of $\RR$ are ``bad'' in that
they don't allow for \autoref{eq:om-finite-additivity} to hold.
In this lecture we will focus on finding a subcollection of $2^\RR$
that consists of ``good'' sets. Such a subcollection should
of course contain $\varnothing$ and $X$, as well as
being closed under complements and countable unions.

\begin{definition}[$\sigma$-algebra]
    \label{def:sigma-algebra}
    A \vocab{$\sigma$-algebra} (or $\sigma$-field) on a set $X$
    is a subcollection $\mathcal S\subseteq 2^X$ such that
    \begin{enumerate}[(1)]
	\setcounter{enumi}{-1}
        \ii $\varnothing \in\mathcal S$.
	\ii if $A\in\mathcal S$, $X\setminus A\in\mathcal S$.
	\ii if $A_1,A_2,\ldots\in\mathcal S$, then
	$\bigcup_{k=1}^\infty A_k\in\mathcal S$.
    \end{enumerate}
    One could also replace (0) with $X\in\mathcal S$,
    or (1) with closure under set difference
    (so $A,B\in\mathcal S$ means $B\setminus A\in\mathcal S$),
    or (2) with closure under countable intersections.
\end{definition}

\begin{example}
    Let $X$ be any set. The following are $\sigma$-algebras on $X$:
    \begin{enumerate}[(1)]
	\ii $\mathcal S = \{\varnothing, X\}$
	\ii $\mathcal S = 2^X$
	\ii $\mathcal S = \{ A\subseteq X : A\text{ is countable or cocountable}\}$
    \end{enumerate}
\end{example}

\begin{lemma}
    \label{lem:intersect-sigma-algebras}
    If $\{\mathcal S_\alpha\}_{\alpha\in A}$ is a collection
    of $\sigma$-algebras on $X$, then
    \begin{equation*}
	\bigcap_{\alpha\in A}\mathcal S_\alpha
    \end{equation*}
    is also a $\sigma$-algebra on $X$.
\end{lemma}

\begin{corollary}
    \label{coro:generated-sigma-algebra}
    If $\mathcal S_0$ is a subset of $2^X$, then
    \begin{equation*}
	\sigma(\mathcal S_0)\coloneq
	\bigcap_{\substack{\text{$\sigma$-algebra }\mathcal S\\
	\mathcal S_0\subseteq\mathcal S}} \mathcal S
    \end{equation*}
    is the minimal $\sigma$-algebra on $X$ containing $\mathcal S_0$;
    this is the $\sigma$-algebra \vocab{generated by $\mathcal S_0$}.
\end{corollary}

The primary example of this is
\begin{definition}[Borel $\sigma$-algebra]
    \label{def:borel-algebra}
    The \vocab{Borel $\sigma$-algebra} on $\RR$ is the
    $\sigma$-algebra generated by the collection of open intervals in $\RR$
    (or, equivalently, the collection of open subsets of $\RR$).
\end{definition}

\begin{prop}
    \label{prop:lebesgue-algebra}
    Let $\mathcal M$ be the set
    \begin{equation*}
	\mathcal M = \{ A\subseteq\RR :
	\text{for all }\eps > 0,\,\text{there exists }
	A_\eps\overset{\text{closed}}\subseteq\RR\text{ s.t. }
	A_\eps\subseteq A\text{ and }\abs{A\setminus A_\eps}<\eps
        \}
    \end{equation*}
    Then:
    \begin{enumerate}[(1)]
        \ii $\mathcal M$ is a $\sigma$-algebra on $\RR$.
	\ii If $A\subseteq\RR$ is closed, then $A\in\mathcal M$.
	\ii If $A\subseteq\RR$ has $\abs A = 0$, then $A\in\mathcal M$.
    \end{enumerate}
\end{prop}

\begin{proof}
    (2) and (3) are obvious (take $A_\eps$ to be $A$ and $\varnothing$ respectively).

    For (1), we need to check the properties of a $\sigma$-algebra.
    First, obviously $\varnothing, X\in\mathcal M$.
    
    Next, we will show that $\mathcal M$ is closed under countable intersections.
    Let $\eps>0$ and let $A_1$, $A_2$, \dots be a sequence of subsets of $\RR$.
    Then, for each $k$, there is a set $A_{k;\eps/2^k}$ such that
    \begin{equation*}
	\abs{A_k\setminus A_{k;\eps/2^k}} < \frac\eps{2^k}.
    \end{equation*}
    Define $A = \bigcap_{k=1}^\infty A_k$ and
    $A_\eps=\bigcap_{k=1}^\infty A_{k;\eps/2^k}$. Then,
    \begin{equation*}
	\abs{A\setminus A_\eps}
	\le \abs{\bigcup_{k=1}^\infty (A_k\setminus A_{k;\eps/2^k})} \le
	\sum_{k=1}^\infty\abs{A_k\setminus A_{k;\eps/2^k}}\le\eps.
    \end{equation*}

    Finally, we show that $\mathcal M$ is closed under complements.
    Letting $\eps > 0$, we start with the case where $A\in\mathcal M$
    and $\abs A < \infty$; this allows us to find an oepn set
    $U_{\eps/2}$ such that $\abs{U_{\eps/2}} < \abs A + \eps/2$,
    using the definition of outer measure.
    We also have a closed set $A_{\eps/2}$ such that
    $A_{\eps/2}\subseteq A$ and $\abs{A\setminus A_{\eps/2}} < \eps/2$
    due to $A$ being in $\mathcal M$. We would like to show that
    \begin{equation*}
	\abs{(\RR\setminus A)\setminus(\RR\setminus U_{\eps/2})}
	= \abs{U_{\eps/2}\setminus A} < \eps
    \end{equation*}
    because $\RR\setminus U_{\eps/2}$ is a closed subset of $\RR$
    contained in the complement of $A$.
    By monotonicity, \autoref{prop:om-open-closed-addend},
    and subadditivity, we have
    \begin{equation*}
	\abs{U_{\eps/2}\setminus A} \le
	\abs{U_{\eps/2}\setminus A_{\eps/2}} =
	\abs{U_{\eps/2}}-\abs{A_{\eps/2}} <
	\abs A + \eps/2 - \abs A + \abs{A\setminus A_{\eps/2}} < \eps
    \end{equation*}
    Now, let $A\in\mathcal M$ be arbitrary.
    Since $\mathcal M$ contains the closed sets $X$ and $[-k,k]$
    for any $k$, as well as being closed under countable intersections,
    we have $A_k = A\cap[-k,k]\in\mathcal M$
    (padding with copies of $X$ to fill
    the finite intersection into a countale one),
    which has bounded outer measure. Then,
    \begin{equation*}
	\RR\setminus A = \RR\setminus\left(\bigcup_{k=1}^\infty A_k\right) =
    \bigcap_{k=1}^\infty(\RR\setminus A_k)
    \end{equation*}
    which by the $\abs A < \infty$ case and closure under countable intersections
    implies $\RR\setminus A\in\mathcal M$.
\end{proof}

This set $\mathcal M$ contains $\mathcal B$
(because it contains all closed, hence all open, subsets of $\RR$),
and it is called the \vocab{Lebesgue $\sigma$-algebra} on $\RR$.

\begin{prop}
    [Criteria for Lebesgue sets]
    \label{prop:lebesgue-criteria}
    Let $A\subseteq\RR$. The following are equivalent:
    \begin{enumerate}[(a)]
        \ii $A\in\mathcal M$, i.e.\ for each $\eps>0$ there exists
	$A^-_\eps\subseteq A$ that is closed in $\RR$
	such that $\abs{A\setminus A_\eps^-} < \eps$.
	\ii There exists a sequence $A_1,A_2,\ldots\subseteq\RR$ of closed sets
	such that each $A_k$ is a subset of $A$ and
	$\abs{A\setminus\bigcup_{k=1}^\infty A_k} = 0$.
	\ii There exists $A^-\in\mathcal B$ such that $A^-\subseteq A$
	and $\abs{A\setminus A^-} = 0$.
	\ii For each $\eps>0$ there exists $A_\eps^+\supseteq A$
	that is open in $\RR$ such that $\abs{A_\eps^+\setminus A}<\eps$.
	\ii There exists a sequence $A_1,A_2,\ldots\subseteq\RR$
	of open sets such that each $A_k$ contains $A$ and
	$\abs{\bigcap_{k=1}^\infty A_k\setminus A} = 0$.
	\ii There exists $A^+\in\mathcal B$ such that
	$A^+\supseteq A$ and $\abs{A^+\setminus A} = 0$.
    \end{enumerate}
\end{prop}

\begin{proof}
    (a)$\implies$(b). Take $A_n$ to be a closed subset of $\RR$ contained in $A$
    such that $\abs{A\setminus A_n} <1/n$. Then
    \begin{equation*}
	\abs{A\setminus\bigcup_{k=1}^\infty} \le \abs{A\setminus A_n}<\frac 1n
    \end{equation*}
    for all $n$, hence the left-hand side is zero.

    (b)$\implies$(c). Take $A^-$ to be the set $\bigcup_{k=1}^\infty$ above.
    
    (c)$\implies$(a). We can write
    \begin{equation*}
        A = A^-\cup(A\setminus A^-).
    \end{equation*}
    Then $A^-\in\mathcal B\subseteq\mathcal M$ and $A\setminus A^-\in\mathcal M$
    because it has measure zero, so $A$ is $\mathcal M$
    because it is the union of two sets in $\mathcal M$.

    (a)$\implies$(d). Since $A\in\mathcal M$, $\RR\setminus A\in\mathcal M$,
    so there exists a closed subset $A_\eps^-\subseteq\RR$
    contained in $\RR\setminus A$ such that
    $\abs{(\RR\setminus A)\setminus A_\eps^-} < \eps$.
    Letting $A_\eps^+ = \RR\setminus A_\eps^-$, we have
    \begin{equation*}
	\abs{A_\eps^+\setminus A}
	= \abs{(\RR\setminus A)\setminus(\RR\setminus A_\eps^+)}
	= \abs{(\RR\setminus A)\setminus A_\eps^-} < \eps.
    \end{equation*}

    (d)$\implies$(e) and (e)$\implies$(f) proceed
    in an analogous fashion to (a)$\implies$(b) and (b)$\implies$(c).

    (f)$\implies$(a). We can write $A = A^+\setminus(A^+\setminus A)$,
    and the two sets in the difference are in $\mathcal M$,
    so $A$ is also in $\mathcal M$.
\end{proof}

It turns out that this set $\mathcal M$ consists of ``good'' subsets of $\RR$:
if $A,B\in\mathcal M$ are disjoint, then we will show next lecture that
\begin{equation*}
    \abs{A\cup B} = \abs A + \abs B.
\end{equation*}

\section{September 4}\label{sec:sep4}
\begin{prop}
    \label{prop:om-lebesgue-addend}
    If $A,B\subseteq\RR$, $A\cap B = \varnothing$, and
    at least one of $A$ and $B$ is in $\mathcal M$, then
    $\abs{A\cup B} = \abs A + \abs B$.
\end{prop}

\begin{proof}
    Suppose $A\in\mathcal M$. Let $\eps>0$ and
    $A_\eps^-\overset{\text{closed}}\subseteq\RR$ such that
    $A_\eps^-\subseteq A$ and $\abs{A\setminus A_\eps^-} < \eps$.
    Then, $A_\eps^-\cap B = \varnothing$ and
    \begin{equation*}
	\abs{A\cup B}\ge \abs{A_\eps\cup B} =
	\abs{A_\eps^-} + \abs B \ge
	(\abs A - \abs{A\setminus A_\eps^-}) + \abs B
	> \abs A + \abs B - \eps
    \end{equation*}
    for all $\eps>0$,
    using \autoref{prop:om-open-closed-addend} for the equality.
\end{proof}

\begin{corollary}
    \label{coro:lebesgue-countable-additivity}
    If $A_1,A_2,\ldots\in\mathcal M$ are pairwise disjoint,
    \begin{equation*}
	\abs{\bigcup_{n=1}^\infty A_n} = \sum_{n=1}^\infty \abs{A_n}.
    \end{equation*}
\end{corollary}

\begin{proof}
    Noting that \autoref{prop:om-lebesgue-addend} holds
    for finitely many pairwise disjoint sets in $\mathcal M$,
    we have that for any positive integer $N$,
    \begin{equation*}
	\abs{\bigcup_{n=1}^\infty A_n} \ge
	\abs{\bigcup_{n=1}^N A_n} =
	\sum_{n=1}^N \abs{A_n}
    \end{equation*}
    so in the limit we get
    \begin{equation*}
	\abs{\bigcup_{n=1}^\infty A_n} = \sum_{n=1}^\infty \abs{A_n}.\qedhere
    \end{equation*}
\end{proof}

\begin{definition}[Measure]
    \label{def:measure}
    Let $\mathcal S\subseteq 2^X$ be a $\sigma$-algebra on a set $X$.
    A \vocab{measure} $\mu\colon\mathcal S\to[0,\infty]$
    is a function such that
    \begin{itemize}
        \ii $\mu(\varnothing) = 0$
	\ii if $A_1,A_2,\ldots\in\mathcal S$ are pairwise disjoint,
	then $\mu(\bigcup_{n=1}^\infty A_n) = \sum_{n=1}^\infty \mu(A_n)$.
    \end{itemize}
    The triple $(X,\mathcal S, \mu)$ is called a \vocab{measure space}.
\end{definition}

For example, $(\RR, \mathcal M, \mu = \abs{\bullet}_{\mathcal M})$
is a measure space, as is $(\RR, \mathcal S, \mu = \abs{\bullet}_{\mathcal S})$
for any sub-$\sigma$-algebra $\mathcal S\subseteq\mathcal M$.
We say that the former is a \vocab{complete measure space}
because if $A\subseteq B\subseteq\RR$ with $B\in\mathcal M$ and $\abs B = 0$,
then $A\in\mathcal M$ as well; it turns out that
$(\RR, \mathcal B, \abs\bullet)$ is not a complete measure space
(will be proven next time).

In this vein, we can take a measure space $(X,\mathcal S_0,\mu_0)$
and consider its \vocab{completion}, which is a measure space
$(X,\mathcal S, \mu)$ such that $\mu|_{\mathcal S_0} = \mu_0$
and $\mathcal S$ is as small as possible. More constructively:

\pagebreak

\begin{prop}
    Suppose $(X,\mathcal S_0,\mu_0)$ is a measure space
    and let
    \begin{equation*}
        \mathcal N = \{A\subseteq X : A\subseteq B\text{ for }B\in\mathcal S_0
	\text{ s.t. }\mu_0(B) = 0\}.
    \end{equation*}
    Then, define $S = \{A\cup N : A\in\mathcal S_0, N\in\mathcal N\}$
    and the function $\mu(A\cup N) \coloneq \mu_0(A)$
    given $A\cup N\in\mathcal S$.
    The tuple $(X,\mathcal S, \mu)$ is a measure space and the unique
    completion of $(X,\mathcal S_0,\mu_0)$.
\end{prop}

\begin{proof}
    A completion of $(X,\mathcal S_0,\mu_0)$ must contain
    $\mathcal S_0$ and $\mathcal N$ by definition,
    and then the definition of $\mu$ is the only plausible way
    to extend $\mu_0$ to $\mathcal S$,
    so it is clear that $(X,\mathcal S, \mu)$ is the unique completion
    of $(X,\mathcal S_0,\mu_0)$ if the former is a measure space.

    Say that $B\in\mathcal S_0$ is a \emph{null set} if $\mu_0(B) = 0$.
    First, $\mathcal S$ is a $\sigma$-algebra because
    \begin{itemize}
        \ii it contains $\varnothing\in\mathcal S_0$;
	\ii if $A\cup N\in\mathcal S$ then, letting $B\in\mathcal S$
	be a set null set containing $N$, we have
	\begin{equation*}
	    (A\cup N)^c\setminus(A\cup B)^c
	    = (A\cup B)\setminus (A\cup N)
	    = B\setminus N
	\end{equation*}
	which is a subset of the null set $B$, hence
	$(A\cup N)^c = (A\cup B)^c \cup (B\setminus N)$
	is in $\mathcal S$;
	\ii if $A_1\cup N_1, A_2\cup N_2,\ldots\in\mathcal S$ then
	\begin{equation*}
	    \bigcup_{n=1}^\infty A_n\cup N_n =
	    \bigcup_{n=1}^\infty A_n\cup\bigcup_{n=1}^\infty N_n
	    \subseteq \bigcup_{n=1}^\infty A_n\cup\bigcup_{n=1}^\infty B_n
	\end{equation*}
	where $B_n$ is a null set containing $N_n$,
	and the union of the $B_n$ is also a null set because
	its measure is bounded above by the sum of the $\mu_0(B_n)$
	by subadditivity, so $\mathcal S$ is closed under
	countable unions.
    \end{itemize}
    Then, to see that $\mu$ is a well-defined measure on $\mathcal S$,
    if a set can be expressed in two ways as $A_1\cup N_1 = A_2\cup N_2$
    where $A_1,A_2\in\mathcal S_0$ and $N_1$, $N_2$ are subsets of null sets
    we must have $\mu_0(A_1) = \mu_0(A_2)$ because
    \begin{equation*}
	A_1\setminus A_2
	= (N_2\setminus A_2)\setminus((N_1\setminus A_2)\setminus(A_1\setminus A_2))
    \end{equation*}
    and the right-hand side being a subset of a null set
    implies $\mu_0(A_1\setminus A_2) = 0$.

    Since $\mu(\varnothing) = \mu_0(\varnothing) = 0$
    it remains to check countable additivity:
    given sets $A_1\cup N_1,A_2\cup N_2,\ldots\in\mathcal S$
    that are pairwise disjoint,
    we can use the fact that $\bigcup_{n=1}^\infty N_n$
    is a subset of a null set to get
    \begin{equation*}
	\mu\left(\bigcup_{n=1}^\infty A_n\cup N_n\right)
	= \mu_0\left(\bigcup_{n=1}^\infty A_n\right)
	= \sum_{n=1}^\infty \mu_0(A_n)
	= \sum_{n=1}^\infty \mu(A_n\cup N_n).\qedhere
    \end{equation*}
\end{proof}

The definition of a measure is fairly minimal;
we now prove that this additivity condition subsumes
many other useful properties that we'd like $\mu$ to satisfy.

\pagebreak

\begin{prop}
    \label{prop:meas-space-props}
    Let $(X,\mathcal S, \mu)$ be a measure space.
    \begin{itemize}
	\ii[(0)] If $A_1,\dots,A_n\in\mathcal S$ are pairwise disjoint,
	then $\mu(\bigcup_{n=1}^N A_n) = \sum_{n=1}^N\mu(A_n)$.
	\ii[(1)] If $A,B\in\mathcal S$ and $A\subseteq B$ then $\mu(A)\le\mu(B)$.
	\ii[(1$'$)] If $\mu(A) < \infty$ then $\mu(A) = \mu(B) - \mu(B\setminus A)$.
	\ii[(2)] If $A_1,A_2,\ldots\in\mathcal S$ then
	$\mu(\bigcup_{n=1}^\infty A_n)\le \sum_{n=1}^\infty \mu(A_n)$.
	\ii[($2^\subset$)]If also $A_1\subseteq A_2\subseteq\cdots$
	then $\mu(\bigcup_{n=1}^\infty A_n) = \lim_{n\to\infty} \mu(A_n)$.
	\ii[($2^\supset$)]If also $A_1\supseteq A_2\supseteq\cdots$
	and $\mu(A_k) < \infty$ for some $k\in\NN$
	then $\mu(\bigcap_{n=1}^\infty A_n) = \lim_{n\to\infty} \mu(A_n)$.
    \end{itemize}
\end{prop}

\begin{proof}
    (0) Just pad the union with copies of $\varnothing$.

    (1) and (1$'$) By (0) we have
    \begin{equation*}
        \mu(A) + \mu(B\setminus A) = \mu(B)
    \end{equation*}
    so $\mu(B\setminus A) \ge 0$ proves (1);
    if $\mu(A)$ is finite then we can subtract it to get (1$'$).

    (2) Let $B_1 = A_1$ and $B_n = A_n\setminus\bigcup_{k=1}^{n-1}A_k$ for $n>1$.
    Then, since $B_n\subseteq A_n$ for all $n$,
    by countable additivity and (1) we get
    \begin{equation*}
	\mu\left(\bigcup_{n=1}^\infty A_n\right) =
	\mu\left(\bigcup_{n=1}^\infty B_n\right) =
	\sum_{n=1}^\infty \mu(B_n) \le \sum_{n=1}^\infty \mu(A_n).
    \end{equation*}

    ($2^\subset$) Define $B_n$ as in (2). Then for any $n\in\NN$,
    \begin{equation*}
	\mu(A_n) =
	\mu\left(\bigcup_{k=1}^n A_k\right) = \mu\left(\bigcup_{k=1}^n B_k\right)
	= \sum_{k=1}^n \mu(B_k),
    \end{equation*}
    so
    \begin{equation*}
	\lim_{n\to\infty}\mu(A_n) = \sum_{k=1}^\infty \mu(B_k)
	= \mu\left(\bigcup_{k=1}^\infty B_k\right)
	= \mu\left(\bigcup_{k=1}^\infty A_k\right).
    \end{equation*}

    ($2^\supset$) We can assume WLOG that $\mu(A_1) < \infty$,
    since if $A_1$, \dots, $A_{k-1}$ aren't finite, then
    we can simply discard them. Now define $B_n = A_1\setminus A_n$.
    This is an increasing sequence of sets, so by ($2^\subset$),
    \begin{equation*}
	\mu\left(A_1\setminus \bigcap_{n=1}^\infty A_n\right) =
	\mu\left(\bigcup_{n=1}^\infty A_1\setminus A_n\right) =
	\lim_{n\to\infty} \mu(A_1\setminus A_n).
    \end{equation*}
    All the sets in the equation above are finite, so by $(1')$
    \begin{equation*}
	\mu(A_1) - \mu\left(\bigcap_{n=1}^\infty A_n\right)
	= \mu(A_1) - \lim_{n\to\infty} \mu(A_n)
    \end{equation*}
    and then we can cancel out the $\mu(A_1)$ because it is finite.
\end{proof}

\section{September 9}\label{sec:sep9}
In order to talk about integration we can't just talk about
any old arbitrary function $f\colon X\to\RR$; we need to restrict
to certain ``good'' functions.

\begin{definition}[Measurable function]
    \label{def:measurable-function}
    Let $(X,\mathcal S)$ and $(Y,\mathcal T)$ be a measurable spaces
    (i.e. sets equipped with $\sigma$-algebras).
    A function $f\colon X\to Y$ is \vocab{$(\mathcal T,\mathcal S)$-measurable}
    (or just \vocab{measurable}) if $f\inv(B)\in\mathcal S$
    for all $B\in\mathcal T$.
\end{definition}

\begin{exercise}
    Show that if $\mathcal S = 2^X$ or $\mathcal T = \{\varnothing,Y\}$,
    then all functions $f\colon X\to Y$ are measurable.
\end{exercise}

\begin{lemma}
    \label{lem:measurable-via-generators}
    If $(X,\mathcal S)$ and $(Y,\mathcal T)$ are measurable spaces
    and $\mathcal T_0\subseteq\mathcal T$ generates $\mathcal T$,
    then $f\colon X\to Y$ is measurable if and only if
    $f\inv(B)\in\mathcal S$ for all $B\in\mathcal T_0$.
\end{lemma}

\begin{proof}
    If $f$ is measurable then the result is trivial;
    the opposite direction is HW02\#1.
\end{proof}

When speaking about functions $f\colon X\to\RR$ or $X\to[-\infty,\infty]$,
we will take for granted that the $\sigma$-algebra attached to the codomain
is the Borel one. For $[-\infty,\infty]$, we topologize it first by
decreeing intervals of the form $[-\infty, a)$ and $(b,\infty]$
to be open for all $a,b\in\RR$; the corresponding Borel $\sigma$-algebra
is the one generated by the open sets in this space. This $\sigma$-algebra
will be denoted $B_\infty$.
(We will shortly see why letting $\RR$ have the Lebesgue $\sigma$-algebra
is not a good idea.)

\begin{definition}[Lebesgue/Borel measurable]
    \label{def:lebesgue-borel-meas-function}
    A function $f\colon\RR\to\RR$ or to $[-\infty,\infty]$
    is \vocab{Lebesgue measurable} (resp. \vocab{Borel measurable})
    if the pre-image of a Borel set is Lebesgue (resp. Borel).
\end{definition}

Of course, by \autoref{lem:measurable-via-generators},
it suffices to check the pre-images of only a generating set,
\begin{exercise}
    Show that the set of rays $(a,\infty)$ generates $\mathcal B$
    and the set of rays $(a,\infty]$ generates $\mathcal B_\infty$.
    Find other ``nice'' generating sets for these $\sigma$-algebras
    (note also that as per their definitions, the $\sigma$-algebras
    are generated by the open sets in their respective spaces).
\end{exercise}

\begin{corollary}
    \label{coro:cont-monotone-meas}
    If $f\inv\colon\RR\to\RR$ is continuous or monotone,
    then $f$ is Borel (hence Lebesgue) measurable.
\end{corollary}

\begin{proof}
    For continuous functions, the pre-image of an open set is
    open hence Borel.

    For monotone functions, assume WLOG that $f\colon\RR\to\RR$ is
    non-decreasing. Then $f\inv((a,\infty))$ either looks like
    $(a',\infty)$ or $[a,\infty)$, both of which are Borel.
\end{proof}

\begin{corollary}
    \label{coro:contrapositive-measurable}
    If $f\colon\RR\to\RR$ is continuous and $B\subseteq\RR$
    is a subset such that $f\inv(B)\subseteq\RR$ is not Lebesgue,
    then $B\subseteq\RR$ is not Borel.
\end{corollary}

\begin{proof}
    This is just the contrapositive of the statement
    ``if $B$ is Borel, then $f\inv(B)$ is Lebesgue''
    that follows due to $f$ being Lebesgue measurable.
\end{proof}

\begin{prop}
    \label{prop:lebesgue-not-borel-set}
    Let $C$ be the Cantor set.
    There exists a continuous function $f\colon\RR\to\RR$
    and $C'\subseteq C$ such that $f\inv(C')$ is not Lebesgue.
\end{prop}

\begin{proof}
    Recall that the Cantor function $\Lambda\colon[0,1]\to[0,1]$
    is continuous, non-decreasing, and constant on
    every interval contained in $[0,1]\setminus C$.

    Define $g\colon\RR\to\RR$ by
    \begin{equation*}
        g(x) = \begin{cases}
	    x & x\le 0\\
	    x + \Lambda(x) & x\in[0,1]\\
	    x + 1 & x\ge 1
        \end{cases}
    \end{equation*}
    \todo{wtf am i doing}
\end{proof}

There are two main takeawaays of this:
\begin{enumerate}[(1)]
    \ii If we used the Lebesgue $\sigma$-algebra instead of the Borel one
    when $\RR$ is the codomain, we would have continuous functions
    that aren't measurable, which is horrific because of course
    continuous functions are integrable.
    \ii There exists $C'\subseteq\RR$ that is Lebesgue but not Borel,
    so $\mathcal B\subset\mathcal M$ is a proper containment.
\end{enumerate}

The real nice thing about measurable functions is that
they behave very well under convergence.
\begin{prop}
    \label{prop:inf-sup-measurable}
    If $(X,\mathcal S)$ is a measurable space and
    $f_1,f_2,\ldots\colon X\to[=\infty,\infty]$ are
    $\mathcal S$-measurable, then the functions
    $\inf f_k$, $\sup f_k$, $\liminf f_k$, $\limsup f_k$
    are all $\mathcal S$-measurable.
    In particular, if $f_1$, $f_2$, \dots converges pointwise
    to $f\colon X\to[-\infty,\infty]$, then $f$ is measurable.
\end{prop}
\todo{define $\mathcal S$-measurable above}

\begin{proof}
    \todo{pf (shouldn't be hard)}
\end{proof}

\section{September 11}\label{sec:sep11}











\end{document}

